\section{上采样与分辨率恢复}
\label{sec:14-Deep-Learning/04-Convolutional-Networks/04}

一个四倍超分模型，PSNR 比基线高了 0.6 dB，验证集指标一路向好。把输出图放到原尺寸盯着看，平坦的天空和墙面上浮着一层规则的方格纹——间距恰好四个像素，横竖都有，像蒙了一张细网。指标不降反升，因为这层纹的幅度很小；可它是规则的、周期的，人眼一眼就能挑出来。

这层纹不是训练不充分，也不是数据里的噪声。它是上采样算子本身的结构缺陷，在网络还没开始训练之前就已经写好了。

\subsection{把尺寸还回去和把信息还回去是两件事}

编码器把 \(512\times512\) 一路降到 \(16\times16\)。分割要逐像素输出，超分要输出比输入更大的图，生成要从一个短向量长出整张图——都得把空间维重新撑开。

最直接的想法是把卷积倒过来。本章第一篇把 stride-1 卷积写成了矩阵 \(y = Cx\)，\(C\) 是一个 Toeplitz 矩阵；带 stride 的卷积对应的 \(C\) 行数少于列数。转置卷积做的事是

\[
x' = C^{\trans} y',
\]

形状对得上：\(C^{\trans}\) 把 \(T_{\mathrm{out}}\) 维送回 \(T_{\mathrm{in}}\) 维。这个算子在框架里的另一个名字是卷积的梯度算子——如果前向是 \(y=Cx\)，那么反向传播里 \(\partial L/\partial x = C^{\trans}(\partial L/\partial y)\)，正是同一个乘法（推导见\hyperref[sec:14-Deep-Learning/02-Losses-and-Backpropagation/02]{``一次完整的前向与反向传播''}）。转置卷积和卷积的反传共用一份实现，这不是巧合，它们本来就是同一个线性映射。

要紧的是 \(C^{\trans}\) 不是 \(C^{-1}\)。带 stride 的 \(C\) 列数多于行数，根本不可逆；\(C^{\trans}C\) 的秩至多是 \(T_{\mathrm{out}}\)，远小于 \(T_{\mathrm{in}}\)。转置卷积的输出被钉在 \(C^{\trans}\) 的像空间里，那是一个维数不超过 \(T_{\mathrm{out}}\) 的子空间。原始高分辨率信号中落在 \(\nullsp(C)\) 里的那部分，前向时被映成零，任何以 \(y\) 为唯一输入的映射都无法把它区分出来。换句话说，尺寸能还回去，信息不能。

这条结论不依赖线性。解码器无论多深多非线性，都只是低分辨率特征的一个函数；同一份低分辨率特征对应的所有高分辨率输入，在它眼里完全一样，输出必然相同。凡是解码器``补''出来的细节，来源只能是训练数据里的先验，而不是这张图本身。超分模型把模糊的车牌变清晰时写出来的数字，正是这么来的。

\subsection{覆盖次数不均匀就是棋盘纹的全部原因}

转置卷积的具体动作换个说法更清楚：输入的每一个位置 \(i\)，把整个核乘上 \(y'_i\) 后，加写到输出的 \(si,\,si+1,\ldots,si+k-1\) 这 \(k\) 个位置上。相邻两个输入的写入区间错开 \(s\) 个位置，区间长度是 \(k\)。当 \(k>s\) 时区间相互重叠，输出的某些位置被写了多次。

数一数输出位置 \(n\) 被写了几次。它接收 \(i\) 的贡献当且仅当 \(0\le n-si\le k-1\)，即

\[
\frac{n-k+1}{s}\ \le\ i\ \le\ \frac{n}{s},
\]

落在这段区间里的整数个数就是覆盖次数 \(c(n)\)。区间长度恒为 \((k-1)/s\)，但其中含几个整数取决于端点的小数部分，也就是只取决于 \(n \bmod s\)。若 \(s \mid k\)，端点的相对位置对每个剩余类都一样，内部所有 \(n\) 的 \(c(n)\) 都等于 \(k/s\)；若 \(s \nmid k\)，\(c(n)\) 在 \(\lfloor k/s\rfloor\) 与 \(\lceil k/s\rceil\) 之间随 \(n\bmod s\) 周期跳动。

\begin{center}
\begin{tikzpicture}[x=0.44cm,y=0.44cm,>=stealth]
\begin{scope}[shift={(0,0)}]
  \foreach \i in {0,1,2,3} {\draw[gray!70] ({2*\i+0.55},3.7) rectangle ({2*\i+1.45},4.5);}
  \foreach \i in {0,1,2,3} {\foreach \o in {0,1,2} {
     \draw[gray!50,thin] ({2*\i+1},3.7) -- ({2*\i+\o+0.45},2.8);}}
  \foreach \m/\c in {0/1,1/1,2/2,3/1,4/2,5/1,6/2,7/1,8/1} {
     \fill[gray!35] (\m,0) rectangle ({\m+0.9},{1.2*\c});
     \draw[gray!75] (\m,0) rectangle ({\m+0.9},{1.2*\c});}
  \draw[gray] (-0.4,0) -- (9.3,0);
  \node[font=\scriptsize] at (4.5,5.4) {输入 4 个位置};
  \node[font=\scriptsize] at (4.5,-1.1) {\(k=3,\ s=2\)：覆盖次数 1 与 2 交替};
\end{scope}
\begin{scope}[shift={(13,0)}]
  \foreach \i in {0,1,2,3} {\draw[gray!70] ({2*\i+1.05},3.7) rectangle ({2*\i+1.95},4.5);}
  \foreach \i in {0,1,2,3} {\foreach \o in {0,1,2,3} {
     \draw[gray!50,thin] ({2*\i+1.5},3.7) -- ({2*\i+\o+0.45},2.8);}}
  \foreach \m/\c in {0/1,1/1,2/2,3/2,4/2,5/2,6/2,7/2,8/1,9/1} {
     \fill[gray!35] (\m,0) rectangle ({\m+0.9},{1.2*\c});
     \draw[gray!75] (\m,0) rectangle ({\m+0.9},{1.2*\c});}
  \draw[gray] (-0.4,0) -- (10.3,0);
  \node[font=\scriptsize] at (5.0,5.4) {输入 4 个位置};
  \node[font=\scriptsize] at (5.0,-1.1) {\(k=4,\ s=2\)：内部覆盖次数恒为 2};
\end{scope}
\end{tikzpicture}
\end{center}

\emph{核宽不能被步长整除时，输出像素接到的写入次数按相位周期跳动——棋盘纹是这条柱状图的直接可视化。}

图左边就是最常见的那组配置。\(k=3,s=2\)，偶数位置被写两次、奇数位置被写一次，输出的幅度天然带一个周期为 2 的调制。二维上两个方向各调制一次，乘起来得到周期 \(2\times2\) 的方格图案。四倍超分若用 \(k=3,s=4\)，周期变成 4——正是开头那张图上量到的间距。图右边把核宽改成 4，写入区间恰好整齐平铺，内部每个位置都被写两次，调制消失。

训练能不能自己修好？部分能，但要顶着架构的约束。补偿需要不同相位的输出乘以不同的系数，而转置卷积对所有输入位置用的是同一组权重——参数共享是这个算子的定义性质，它没有能力给不同相位分配不同幅度。补偿只能推给后续层，而后续层同样是共享权重的卷积。更麻烦的是，反向传播时梯度沿同样的路径回流，覆盖多的位置也拿到更大的梯度，不均匀在训练动力学里被持续地重新注入。

按 \(s \mid k\) 选核宽能消掉这个周期，代价是核宽被绑定。更常用的两条路绕开了它。一是插值加卷积：先用最近邻或双线性把特征图放大到目标尺寸——这是一个覆盖天然均匀的固定算子——再接 stride 1 的普通卷积去学内容。均匀性由构造保证，与学到什么权重无关；代价是插值自带平滑偏置，且卷积要在放大后的高分辨率上算，开销更大。二是子像素卷积：在低分辨率上一次算出 \(s^2\) 倍的通道，再把通道维重排进空间维。重排是一一对应，不存在重叠写入，但每个相位的输出来自不同的卷积核，若随机初始化，各相位的起点统计量不同，同样会显出纹路；配套的 ICNR 初始化让同一像素块内的 \(s^2\) 个核初始相同，把起点对齐。

三条路的共同点是：先让上采样算子对所有相位一视同仁，再把内容交给可学习的部分。

\subsection{跳跃连接搬的是解码器造不出来的东西}

前面那条零空间论证已经说清了跳跃连接为什么不是可选项。解码器的输入是低分辨率特征，它决定不了的那部分内容，加多少层都决定不了。要恢复边界的确切位置、细纹理的走向、小目标的存在与否，唯一的办法是从还没被抽稀的地方把这些量搬过来。

U-Net 把编码器第 \(l\) 级的特征图直接接到解码器同分辨率的那一级，拼在通道维上。解码器在每一级同时拿到两份东西：从更深处上来的、语义强但空间糊的特征，以及从旁路来的、语义弱但空间准的特征。前者说``这一块是血管''，后者说``边界在这条线上''，两者在同一分辨率上对齐后由卷积融合。分割掩码的边界精度几乎全部来自后者。

边界也要划出来。跳跃连接搬过来的是未经加工的低级特征，编码器早期的噪声、伪影和采集设备的特性一并进来；旁路太宽或太浅时，解码器会发现直接抄浅层比用深层语义更省力，训练损失下降但语义判断退化，表现为大面积区域的类别错误配上极其锐利的边界。另一方面，跳跃连接只能搬来编码器算过的东西——如果高频信息在输入端就被镜头模糊掉了，编码器第一层也没有，旁路搬来的仍是空的。跳跃连接解决的是``网络内部丢失''，解决不了``数据里本来就没有''。

到这里，卷积网络的整套机器已经装齐：局部核把平移结构写进算子，深度和空洞把视野推远，下采样拿分辨率换不变性，上采样加旁路再把分辨率换回来。每一件都是围着同一组先验展开的工程。下一篇退后一步，问这组先验本身值多少钱——它在什么条件下是巨大的优势，又在什么条件下变成天花板。
